\begin{textsample}{POS Dim 1 – vem\_ed\_07 – Score 2.00 – vem\_ed\_07/t028.txt}  \label{ex:f1_pos_185}
Incrementar a \textbf{interação} BOAS PRÁTICAS Camila Kami leciona inglês e trabalha com PCIs desde o segundo semestre de 2019, quando \textbf{iniciou} um projeto sobre equipamentos médico-hospitalares com Luiz Roberto Madureira Iório, professor do curso de Sistemas Biomédicos (Fatec Bauru), e Liesa Kyer, da área de enfermagem (BridgeValley Community \& Technical College, EUA).
Houve ainda uma segunda edição desse PCI em 2020.
Neste primeiro semestre de 2021, Camila atuou em dois projetos: um sobre planejamento estratégico de marketing, com os professores Antonio Cesar Dall’Evedove (Fatec Garça) e Paul Barretta (Wagner College, EUA).
E outro sobre plano de negócios, com Kleber Luiz Nardoto Milaneze (Fatec Bauru) e Tiffany Lean Macquarrie (PennState Beaver, EUA).
A professora compartilha sua experiência a \textbf{seguir}: O planejamento \textbf{é} fundamental para traçar objetivos de aprendizagem e delinear atividades \textbf{significativas} para os alunos de ambas as instituições.
Várias reuniões com o parceiro internacional \textbf{são} necessárias para amadurecer o projeto e solucionar dúvidas. \textbf{É} muito importante sugerir aos alunos uma lista de \textbf{questões} para a apresentação individual ou em grupo da primeira fase do projeto, o icebreaker [ quebra-gelo ].
Um vídeo com várias informações sobre os alunos facilita a \textbf{interação} e as identificações entre os parceiros.
Além disso, \textbf{é} interessante sugerir uma ferramenta para a gravação dos vídeos.
O PowerPoint narrado apresenta-se como uma ótima sugestão devido à possibilidade de narrar \textbf{cada} slide separadamente e à praticidade em apagar apenas as gravações \textbf{que} não ficaram boas.
O arquivo pode \textbf{ser} salvo no formato MP4.
Para \textbf{que} o aluno se sinta mais seguro, \textbf{é} recomendável \textbf{que} o professor de inglês revise o texto a \textbf{ser} narrado, \textbf{identificando} trechos \textbf{que} talvez não \textbf{fossem} compreendidos pelo parceiro internacional.
Ferramentas como SpeechNinja e YouGlish auxiliam com a pronúncia.
Incrementar a \textbf{interação} entre os alunos \textbf{é} um grande desafio.
Como na maioria das vezes o projeto \textbf{é} desenvolvido de \textbf{forma} assíncrona, \textbf{é} muito importante agendar, pelo menos, uma reunião síncrona com \textbf{todos} os participantes logo no \textbf{início} do projeto, uma oportunidade para \textbf{que} os alunos vejam os colegas.
Nessa reunião, os professores de ambas as instituições mediam a comunicação entre os alunos, encorajando-os a participar.
A avaliação do projeto \textbf{é} facilitada por meio da criação de rubricas: os professores explicitam os critérios a \textbf{serem} avaliados.
O professor \textbf{que} deseja participar de um PCI precisa estar disposto a conhecer, testar e usar diferentes ferramentas tecnológicas, estar aberto a \textbf{diálogos} constantes com os parceiros internacionais, verificar a entrega das atividades nos prazos definidos, contactar os alunos caso não cumpram alguma etapa, auxiliá-los em suas dúvidas e avaliar o projeto.

% matched lemmas: cada, diálogo, forma, identificar, iniciar, interação, início, que, questão, seguir, ser, significativo, todo
\end{textsample}
